\chapter{Étude de l'existant et méthodologie de travail}
\label{chap:etude_existant}

\section{Introduction}

Avant de concevoir NetSentry, il est utile d'examiner les solutions déjà disponibles
pour la découverte et la supervision des équipements réseau, afin de justifier les choix
retenus. Ce chapitre présente cette étude de l'existant, une analyse SWOT du projet, le
cahier des charges qui a cadré le stage, ainsi que les exigences fonctionnelles et non
fonctionnelles retenues pour NetSentry.

\section{Solutions existantes}

Plusieurs catégories d'outils permettent, à des degrés divers, de découvrir ou de
superviser les équipements d'un réseau local :

\subsection{Outils de scan en ligne de commande}
\begin{itemize}
    \item \textbf{Nmap} : scanner réseau open source, extrêmement puissant et précis
    (découverte d'hôtes, détection de ports et de services), mais qui reste un outil en
    ligne de commande : il ne conserve aucun historique entre deux exécutions et ne
    propose ni tableau de bord, ni mécanisme d'alerte natif.
\end{itemize}

\subsection{Applications grand public}
\begin{itemize}
    \item \textbf{Fing} : application mobile de scan réseau simple d'utilisation,
    orientée particuliers, avec des fonctionnalités d'historique et d'alerte limitées à
    la version payante, et peu adaptée à un déploiement interne sur un serveur d'hôtel.
    \item \textbf{Advanced IP Scanner} : utilitaire Windows léger de découverte
    d'équipements, sans persistance des résultats ni planification automatique.
\end{itemize}

\subsection{Solutions de supervision professionnelles}
\begin{itemize}
    \item \textbf{PRTG Network Monitor} : solution de supervision réseau complète et
    reconnue, mais commerciale, dont le coût de licence et la richesse fonctionnelle
    dépassent largement les besoins d'un simple inventaire d'équipements pour un
    département IT hôtelier.
    \item \textbf{Zabbix / Nagios} : plateformes de supervision open source très
    puissantes, mais dont l'installation, la configuration et la prise en main
    demandent une expertise et un temps d'appropriation importants -- peu compatibles
    avec la contrainte de simplicité d'utilisation pour un personnel non spécialisé en
    réseau.
\end{itemize}

\section{Limites observées}

Cette étude fait ressortir un constat commun : les outils de scan bruts (Nmap) n'offrent
aucune persistance ni interface, tandis que les solutions grand public manquent de
fonctionnalités d'historisation et d'alerte adaptées, et que les plateformes de
supervision professionnelles sont surdimensionnées, coûteuses ou trop complexes pour un
déploiement léger et autonome au sein d'un seul département IT.

\begin{itemize}
    \item Aucune solution ne combine à la fois simplicité de déploiement,
    auto-hébergement et historique exploitable pour ce périmètre restreint ;
    \item Les solutions commerciales impliquent un coût de licence récurrent, non
    justifié pour un besoin d'inventaire d'équipements ;
    \item Les solutions de supervision complètes nécessitent une expertise réseau que le
    personnel IT non spécialisé ne possède pas nécessairement ;
    \item Peu de solutions permettent un déploiement conteneurisé simple, cohérent avec
    les contraintes d'infrastructure de l'hôtel.
\end{itemize}

\section{Analyse SWOT}

Le tableau~\ref{tab:swot} synthétise les forces, faiblesses, opportunités et menaces
identifiées pour le projet NetSentry.

\begin{table}[H]
\centering
\begin{tabular}{|p{6.3cm}|p{6.3cm}|}
\hline
\textbf{Forces} & \textbf{Faiblesses} \\
\hline
Développement sur mesure, adapté au besoin réel du département &
Projet réalisé par un seul stagiaire, sur une durée de deux mois \\
Basé sur Nmap, un outil de scan éprouvé et fiable &
Pas de détection de type d'équipement par apprentissage automatique \\
Auto-hébergé et conteneurisé (Docker), sans coût de licence &
OS detection (\texttt{nmap -O}) exclue du périmètre (nécessite des privilèges root
supplémentaires) \\
Interface simple, pensée pour du personnel non spécialisé en réseau &
Historique GitHub encore synthétique à ce stade du projet \\
\hline
\textbf{Opportunités} & \textbf{Menaces} \\
\hline
Sensibilisation croissante des hôtels à la cybersécurité réseau &
Évolution des équipements vers des adresses MAC aléatoires (confidentialité), limitant
l'identification par fabricant \\
Réutilisable pour d'autres hôtels du groupe ou d'autres PME &
Dépendance à un prestataire externe pour les systèmes métiers, hors périmètre du projet \\
Base pour des évolutions futures (notifications e-mail, classification par type
d'appareil) &
Réseau hôtelier partiellement hors du contrôle direct du département IT interne \\
\hline
\end{tabular}
\caption{Analyse SWOT du projet NetSentry}
\label{tab:swot}
\end{table}

\section{Cahier des charges}

Le cadrage du projet a été formalisé dans un cahier des charges remis en début de
stage, définissant le contexte, les objectifs, le périmètre fonctionnel et non
fonctionnel, ainsi que les contraintes et livrables attendus.

\subsection{Livrables attendus}

\begin{itemize}
    \item Code source de l'application (backend, frontend, scripts de scan) ;
    \item Base de données structurée contenant l'historique des équipements détectés ;
    \item Tableau de bord fonctionnel de visualisation ;
    \item Image(s) Docker et instructions de déploiement ;
    \item Documentation technique et rapport de stage détaillant la démarche, les choix
    techniques et les résultats obtenus.
\end{itemize}

\subsection{Contraintes}

Le projet devait être réalisé dans un délai de deux mois, avec un accès limité aux
systèmes métiers de l'hôtel (gérés par un prestataire externe). Le scan réseau ne
pouvait être exécuté que sur le réseau interne de l'hôtel, avec l'autorisation du
département IT, et la solution devait privilégier une base de données relationnelle
(SQL) en cohérence avec les outils disponibles côté hôtel.

\section{Exigences fonctionnelles}

Les exigences fonctionnelles traduisent les besoins identifiés en fonctionnalités
concrètes pour NetSentry :

\begin{itemize}
    \item \textbf{Découverte des équipements} : scanner le réseau local afin
    d'identifier tous les équipements actuellement connectés (adresse IP, adresse MAC).
    \item \textbf{Identification des équipements} : déterminer, lorsque c'est possible,
    le fabricant (via correspondance MAC/OUI) et les ports/services ouverts.
    \item \textbf{Historique des équipements} : enregistrer chaque scan en base de
    données afin de conserver un historique (nouveaux équipements, équipements connus,
    équipements disparus).
    \item \textbf{Tableau de bord} : visualiser en temps réel la liste des équipements
    connectés, leur statut, et l'historique d'activité du réseau.
    \item \textbf{Détection d'anomalies} : identifier et signaler automatiquement tout
    équipement inconnu apparaissant pour la première fois sur le réseau.
    \item \textbf{Alertes} : notifier l'administrateur IT en cas de détection d'un
    équipement non reconnu, et lui permettre de l'acquitter ou de le classer comme
    connu.
    \item \textbf{Planification} : exécuter les scans automatiquement selon une
    fréquence configurable, sans intervention manuelle systématique.
    \item \textbf{Déploiement} : conteneuriser l'application avec Docker afin de
    permettre un déploiement simple sur un serveur de l'hôtel.
\end{itemize}

\section{Exigences non fonctionnelles}

Les exigences non fonctionnelles garantissent la qualité, la sécurité et la performance
du système, indépendamment de ses fonctionnalités métier :

\begin{itemize}
    \item \textbf{Ergonomie} : le tableau de bord doit rester simple et lisible pour un
    personnel non spécialisé en réseau.
    \item \textbf{Sécurité} : l'accès au tableau de bord doit être protégé par
    authentification (JWT), et les mots de passe stockés de façon chiffrée (hachage
    bcrypt).
    \item \textbf{Confidentialité} : aucune donnée du réseau interne ne doit être
    transmise à des services externes -- l'identification des fabricants s'appuie sur une
    base OUI locale, sans appel API externe.
    \item \textbf{Autonomie} : les scans doivent pouvoir s'exécuter de façon planifiée,
    sans intervention manuelle systématique.
    \item \textbf{Portabilité} : l'application doit être conteneurisée (Docker) afin de
    pouvoir être déployée sur différents environnements.
    \item \textbf{Performance} : un scan complet du réseau local doit s'exécuter en un
    temps raisonnable, sans perturber le trafic réseau existant (usage d'un délai
    d'expiration par hôte).
    \item \textbf{Maintenabilité} : le code doit être structuré (architecture claire
    backend/frontend), documenté et versionné (Git).
\end{itemize}

\section{Méthodologie de travail}

Le projet a été réalisé selon une approche itérative : mise à niveau sur les
fondamentaux réseau, puis développement progressif du backend, du frontend, du module de
détection d'anomalies, et enfin conteneurisation et documentation -- chaque étape étant
validée avant de passer à la suivante. L'organisation du travail au fil du stage (planning
prévisionnel et suivi des tâches) est détaillée au chapitre~\ref{chap:archi_tech}.

\section{Conclusion}

Cette étude de l'existant a permis de justifier le choix d'une solution développée sur
mesure plutôt que l'adoption d'un outil tiers, et l'analyse SWOT ainsi que le cahier des
charges ont permis de cadrer précisément le périmètre du projet. Le chapitre suivant
présente l'analyse et la conception du système NetSentry, à travers ses diagrammes UML.
